% =========================================================
% Systems of Sets
% =========================================================

\section{Systems of Sets}

% ---------------------------------------------------------
% TOOLKIT
% ---------------------------------------------------------
\begin{exposition}
\vspace{0.5em}
\end{exposition}

\begin{tcolorbox}[colback=gray!6, colframe=gray!40, arc=2pt,
  left=6pt, right=6pt, top=4pt, bottom=4pt,
  title={\small\textbf{Comparison of Set Systems}},
  fonttitle=\small\bfseries]
\small
\begin{tabular}{l c c c c c}
\toprule
\textbf{Property} & \textbf{Semiring} & \textbf{Ring} & \textbf{Algebra} & \textbf{$\sigma$-algebra} & \textbf{Topology} \\
\midrule
Contains $\varnothing$            & \checkmark & \checkmark & \checkmark & \checkmark & \checkmark \\
Contains $X$                      & ---        & ---        & \checkmark & \checkmark & \checkmark \\
Closed under $\cap$ (finite)      & \checkmark & \checkmark & \checkmark & \checkmark & \checkmark \\
Closed under $\cup$ (finite)      & ---        & \checkmark & \checkmark & \checkmark & \checkmark \\
Closed under $\cup$ (countable)   & ---        & ---        & ---        & \checkmark & --- \\
Closed under $\cup$ (arbitrary)   & ---        & ---        & ---        & ---        & \checkmark \\
Closed under complement           & ---        & ---        & \checkmark & \checkmark & --- \\
Difference $A \setminus B$        & finite disj.\ union & \checkmark & \checkmark & \checkmark & --- \\
\bottomrule
\end{tabular}
\end{tcolorbox}

\begin{exposition}
\vspace{1em}
\end{exposition}

\begin{exposition}
The definitions and theorems of the previous sections have been about
individual sets and their operations. A second level of structure emerges
when one asks: which \emph{collections} of sets are closed under given
operations? This question is not abstract generality for its own sake. The
three subjects that depend most heavily on this chapter --- topology, measure
theory, and functional analysis --- each begin by specifying a particular type
of collection on a given set $X$. A topology specifies which subsets are
``open.'' A $\sigma$-algebra specifies which subsets are ``measurable.'' In
both cases, the definition is precisely a list of closure conditions: which
set operations applied to members of the collection must produce another
member.
\end{exposition}

\begin{exposition}
Mastering these four structures here, as purely set-theoretic objects, means
that when topology and measure theory are encountered, all the foundational
work is already done. The task will be to understand the \emph{geometric} or
\emph{analytic} meaning of the axioms, not to struggle with the set-theoretic
mechanics.
\end{exposition}

% ---------------------------------------------------------
% Semiring
% ---------------------------------------------------------

\begin{definitionbox}{Definition (Semiring of Sets)}
\begin{definition}[Semiring of Sets]\label{def:semiring-of-sets}
A collection $\mathcal{S}$ of subsets of a set $X$ is a \emph{semiring} if:
\begin{enumerate}[label=(\roman*)]
\item $\varnothing \in \mathcal{S}$;
\item $A, B \in \mathcal{S}$ implies $A \cap B \in \mathcal{S}$;
\item for any $A, B \in \mathcal{S}$, the difference $A \setminus B$ is a
  \emph{finite disjoint union} of members of $\mathcal{S}$: there exist
  pairwise disjoint $C_1, \dots, C_n \in \mathcal{S}$ with
  $A \setminus B = C_1 \sqcup \cdots \sqcup C_n$.
\end{enumerate}
\end{definition}
\end{definitionbox}

\begin{dependencies}
\begin{itemize}
  \item \hyperref[def:subset]{Subset}
  \item \hyperref[def:empty-set]{Empty Set}
  \item \hyperref[def:intersection]{Intersection}
  \item \hyperref[def:set-difference]{Set Difference}
  \item \hyperref[def:union]{Union}
  \item \hyperref[def:pairwise-disjoint]{Pairwise Disjoint Family}
  \item \hyperref[def:finite-infinite]{Finite and Infinite Sets}
\end{itemize}
\end{dependencies}

\begin{example*}[Half-open intervals form a semiring]
\label{ex:semiring-intervals}
Let $X = \mathbb{R}$ and let $\mathcal{S}$ be the collection of all half-open
intervals $[a, b) = \{x \in \mathbb{R} : a \leq x < b\}$, together with
$\varnothing$. Then $\mathcal{S}$ is a semiring:
\begin{itemize}
\item $[a,b) \cap [c,d) = [\max(a,c),\, \min(b,d))$, which is again a half-open
  interval (or $\varnothing$);
\item $[a,b) \setminus [c,d)$ is either $\varnothing$, $[a,c)$, $[d,b)$, or
  $[a,c) \sqcup [d,b)$ --- in every case, a finite disjoint union of half-open
  intervals.
\end{itemize}
This semiring is the natural starting point for constructing Lebesgue measure:
one first defines length on $\mathcal{S}$ (as $b - a$ for $[a,b)$), then
extends to the generated ring and $\sigma$-algebra.
\end{example*}

\begin{remark*}[Why semirings appear in measure theory]
Lebesgue measure is most naturally defined first on a semiring, because
semirings have enough structure to state the countable additivity condition for
a \emph{premeasure}, yet are simple enough that explicit formulas for sets like
$[a,b)$ are easy to write down. The general extension theorem (Carath\'{e}odory's
theorem) then extends the premeasure to the full $\sigma$-algebra. This is why
Kolmogorov \& Fomin introduce semirings before algebras: they are the base
case of the extension hierarchy.
\end{remark*}

% ---------------------------------------------------------
% Ring of sets
% ---------------------------------------------------------

\begin{definitionbox}{Definition (Ring of Sets)}
\begin{definition}[Ring of Sets]\label{def:ring-of-setsxxx}
A collection $\mathcal{R}$ of subsets of $X$ is a \emph{ring of sets} if:
\begin{enumerate}[label=(\roman*)]
\item $\varnothing \in \mathcal{R}$;
\item $A, B \in \mathcal{R}$ implies $A \cup B \in \mathcal{R}$;
\item $A, B \in \mathcal{R}$ implies $A \setminus B \in \mathcal{R}$.
\end{enumerate}
\end{definition}
\end{definitionbox}

\begin{dependencies}
\begin{itemize}
  \item \hyperref[def:subset]{Subset}
  \item \hyperref[def:empty-set]{Empty Set}
  \item \hyperref[def:union]{Union}
  \item \hyperref[def:set-difference]{Set Difference}
\end{itemize}
\end{dependencies}

\begin{remark*}[Derived closure properties of a ring]
A ring is closed under finitely many of each operation. Specifically, if
$\mathcal{R}$ is a ring then:
\begin{itemize}
\item $A \cap B = A \setminus (A \setminus B) \in \mathcal{R}$ (intersection
  follows from set difference);
\item $A \triangle B = (A \setminus B) \cup (B \setminus A) \in \mathcal{R}$
  (symmetric difference follows from union and difference).
\end{itemize}
A ring is closed under all finite Boolean combinations of its elements.
\end{remark*}

\begin{remark*}[Why ``ring'']
The name comes from algebra. Under the operations of symmetric difference
$\triangle$ (playing the role of addition) and intersection $\cap$
(multiplication), a ring of sets is a commutative ring in the algebraic
sense: $\triangle$ is commutative and associative with identity $\varnothing$,
every element is its own inverse ($A \triangle A = \varnothing$), and $\cap$
distributes over $\triangle$.
\end{remark*}

% ---------------------------------------------------------
% Algebra of sets
% ---------------------------------------------------------

\begin{definitionbox}{Definition (Algebra of Sets)}
\begin{definition}[Algebra of Sets]\label{def:algebra-of-sets}
A collection $\mathcal{A}$ of subsets of $X$ is an \emph{algebra of sets}
(or \emph{Boolean algebra of sets}, or \emph{field of sets}) if:
\begin{enumerate}[label=(\roman*)]
\item $X \in \mathcal{A}$;
\item $A \in \mathcal{A}$ implies $A^c = X \setminus A \in \mathcal{A}$;
\item $A, B \in \mathcal{A}$ implies $A \cup B \in \mathcal{A}$.
\end{enumerate}
\end{definition}
\end{definitionbox}

\begin{dependencies}
\begin{itemize}
  \item \hyperref[def:subset]{Subset}
  \item \hyperref[def:complement]{Complement}
  \item \hyperref[def:union]{Union}
\end{itemize}
\end{dependencies}

\begin{remark*}[Derived properties of an algebra]
An algebra is a ring that also contains $X$. Since $\varnothing = X^c$ and
$X \in \mathcal{A}$, every algebra contains $\varnothing$. Closure under
finite intersection follows: $A \cap B = (A^c \cup B^c)^c \in \mathcal{A}$
by De Morgan. An algebra is therefore closed under all finite Boolean
operations: finite unions, finite intersections, complements, and differences.
\end{remark*}

% ---------------------------------------------------------
% Sigma-algebra
% ---------------------------------------------------------

\begin{definitionbox}{Definition ($\sigma$-Algebra)}
\begin{definition}[$\sigma$-Algebra]\label{def:sigma-algebra}
A collection $\mathcal{M}$ of subsets of $X$ is a \emph{$\sigma$-algebra}
(or \emph{$\sigma$-field}) on $X$ if:
\begin{enumerate}[label=(\roman*)]
\item $X \in \mathcal{M}$;
\item $A \in \mathcal{M}$ implies $A^c \in \mathcal{M}$;
\item if $\{A_n\}_{n \in \mathbb{N}}$ is a countable family in $\mathcal{M}$,
  then $\bigcup_{n=1}^{\infty} A_n \in \mathcal{M}$.
\end{enumerate}
The pair $(X, \mathcal{M})$ is called a \emph{measurable space}.
\end{definition}
\end{definitionbox}

\begin{dependencies}
\begin{itemize}
  \item \hyperref[def:subset]{Subset}
  \item \hyperref[def:empty-set]{Empty Set}
  \item \hyperref[def:complement]{Complement}
  \item \hyperref[def:union]{Union}
  \item \hyperref[def:countable]{Countable Set}
\end{itemize}
\end{dependencies}

\begin{remark*}[Derived closure properties of a $\sigma$-algebra]
Every $\sigma$-algebra is an algebra (take the union of a finite family by
padding with $\varnothing$). But a $\sigma$-algebra is additionally closed
under countable intersections: by De Morgan,
\[
\bigcap_{n=1}^{\infty} A_n
\;=\;
\Bigl(\bigcup_{n=1}^{\infty} A_n^c\Bigr)^c
\;\in\; \mathcal{M}.
\]
The $\sigma$ in ``$\sigma$-algebra'' is notation for ``countable'' (from the
German \emph{Summe}), signalling the step up from finite to countable closure.
\end{remark*}

\begin{remark*}[Why countable and not arbitrary]
A $\sigma$-algebra is closed under countable unions but not arbitrary ones.
This is not a deficiency --- it is the right condition for measure theory.
Measure is defined to be countably additive: if $\{A_n\}$ is a countable
disjoint family, then $\mu\bigl(\bigcup A_n\bigr) = \sum \mu(A_n)$. Extending
to uncountable unions would force every set to have either measure zero or
infinite measure (a useless theory). The countability restriction is what makes
the theory nontrivial.

By contrast, a topology (defined below) is closed under \emph{arbitrary}
unions but only \emph{finite} intersections. This is a different trade-off:
open sets are defined by local conditions, and any union of open sets is open
because being open is a pointwise property. Countable intersections of open
sets are \emph{not} required to be open; those are a separate class called
$G_\delta$ sets.
\end{remark*}

% ---------------------------------------------------------
% Topology
% ---------------------------------------------------------

\begin{definitionbox}{Definition (Topology)}
\begin{definition}[Topology]\label{def:topology-defn}
A \emph{topology} on a set $X$ is a collection $\tau$ of subsets of $X$
satisfying:
\begin{enumerate}[label=(\roman*)]
\item $\varnothing \in \tau$ and $X \in \tau$;
\item if $\{U_\alpha\}_{\alpha \in I}$ is any family of sets in $\tau$ (indexed
  by any set $I$), then $\bigcup_{\alpha \in I} U_\alpha \in \tau$;
\item if $U_1, \dots, U_n \in \tau$ (finitely many), then $U_1 \cap \cdots
  \cap U_n \in \tau$.
\end{enumerate}
The pair $(X, \tau)$ is a \emph{topological space}. Members of $\tau$ are
called \emph{open sets}.
\end{definition}
\end{definitionbox}

\begin{dependencies}
\begin{itemize}
  \item \hyperref[def:subset]{Subset}
  \item \hyperref[def:empty-set]{Empty Set}
  \item \hyperref[def:indexed-family]{Indexed Family of Sets}
  \item \hyperref[def:indexed-union]{Indexed Union}
  \item \hyperref[def:indexed-intersection]{Indexed Intersection}
  \item \hyperref[def:finite-infinite]{Finite and Infinite Sets}
\end{itemize}
\end{dependencies}

\begin{remark*}[Why arbitrary unions but only finite intersections]
A topology formalizes the idea of ``nearness'' without a distance function.
The condition that arbitrary unions of open sets are open reflects the
following: if a point $x$ is near some set, and that set is a union of smaller
sets, then $x$ is near one of those smaller sets. This is a purely local
condition and works for any union.

Finite intersections must be open because: if $x$ lies in the interior of
$U_1$ and also in the interior of $U_2$, it lies in the interior of $U_1 \cap
U_2$. This argument works for finitely many sets (take the smallest
neighborhood). For a countably infinite intersection, the ``smallest
neighborhood'' might shrink to nothing --- consider $\bigcap_{n=1}^\infty
(-1/n, 1/n) = \{0\}$ in $\mathbb{R}$, which is not open. So infinite
intersections of open sets are not required to be open.
\end{remark*}

% ---------------------------------------------------------
% Generated structures
% ---------------------------------------------------------

\begin{theorembox}{Theorem (Intersection of $\sigma$-algebras is a $\sigma$-algebra)}
\begin{theorem}[Intersection of $\sigma$-algebras is a $\sigma$-algebra]
\label{thm:sigma-intersection}
Let $\{\mathcal{M}_\alpha\}_{\alpha \in I}$ be any family of $\sigma$-algebras
on $X$ (indexed by any set $I$). Then
\[
\mathcal{M} \;:=\; \bigcap_{\alpha \in I} \mathcal{M}_\alpha
\]
is a $\sigma$-algebra on $X$.
\hyperref[prf:sigma-intersection]{Go to proof.}
\end{theorem}
\end{theorembox}

\begin{remark*}[Standard quantified statement]
$\bigl(\forall\alpha\in I,\; \mathcal{M}_\alpha \text{ is a $\sigma$-algebra on }X\bigr) \;\to\; \bigcap_{\alpha\in I}\mathcal{M}_\alpha \text{ is a $\sigma$-algebra on }X$.\\
Explicitly: $X\in\bigcap\mathcal{M}_\alpha$; $A\in\bigcap\mathcal{M}_\alpha \to A^c\in\bigcap\mathcal{M}_\alpha$; $\forall n,\; A_n\in\bigcap\mathcal{M}_\alpha \to \bigcup_n A_n\in\bigcap\mathcal{M}_\alpha$.
\end{remark*}

\begin{dependencies}
\begin{itemize}
  \item \hyperref[def:sigma-algebra]{$\sigma$-Algebra}
  \item \hyperref[def:indexed-family]{Indexed Family of Sets}
  \item \hyperref[def:intersection]{Intersection}
\end{itemize}
\end{dependencies}

\begin{remark*}[Proof]
Check each axiom: (i) $X \in \mathcal{M}_\alpha$ for all $\alpha$, so
$X \in \bigcap_\alpha \mathcal{M}_\alpha$. (ii) If $A \in \mathcal{M}$ then
$A \in \mathcal{M}_\alpha$ for all $\alpha$, so $A^c \in \mathcal{M}_\alpha$
for all $\alpha$, so $A^c \in \mathcal{M}$. (iii) If $\{A_n\} \subseteq
\mathcal{M}$ then $\{A_n\} \subseteq \mathcal{M}_\alpha$ for all $\alpha$, so
$\bigcup A_n \in \mathcal{M}_\alpha$ for all $\alpha$, so $\bigcup A_n \in
\mathcal{M}$. All three conditions are verified. The same argument works for
algebras, rings, and topologies (with appropriate closure conditions).
\end{remark*}

\begin{definitionbox}{Definition (Generated $\sigma$-Algebra)}
\begin{definition}[Generated $\sigma$-Algebra]\label{def:generated-sigma}
Let $\mathcal{F}$ be any collection of subsets of $X$. The
\emph{$\sigma$-algebra generated by $\mathcal{F}$}, written $\sigma(\mathcal{F})$,
is the smallest $\sigma$-algebra on $X$ that contains $\mathcal{F}$:
\[
\sigma(\mathcal{F})
\;:=\;
\bigcap\,\bigl\{\, \mathcal{M} \;\bigm|\;
\mathcal{M} \text{ is a $\sigma$-algebra on } X
\text{ and } \mathcal{F} \subseteq \mathcal{M} \,\bigr\}.
\]
The collection $\mathcal{F}$ is called a \emph{generating family} for
$\sigma(\mathcal{F})$.
\end{definition}
\end{definitionbox}

\begin{dependencies}
\begin{itemize}
  \item \hyperref[def:sigma-algebra]{$\sigma$-Algebra}
  \item \hyperref[def:subset]{Subset}
  \item \hyperref[def:intersection]{Intersection}
  \item \hyperref[thm:sigma-intersection]{Intersection of $\sigma$-Algebras}
\end{itemize}
\end{dependencies}

\begin{remark*}[Why the definition is valid]
The set of $\sigma$-algebras on $X$ containing $\mathcal{F}$ is nonempty:
$\mathcal{P}(X)$ is a $\sigma$-algebra containing every subset of $X$.
Theorem~\ref{thm:sigma-intersection} guarantees their intersection is again a
$\sigma$-algebra. So $\sigma(\mathcal{F})$ is always well-defined. The same
construction works for generated algebras, generated rings, and (with a small
modification for topologies) generated topologies.
\end{remark*}

\begin{remark*}[How to think about generated structures]
$\sigma(\mathcal{F})$ contains $\mathcal{F}$, and then contains everything
that must be present in any $\sigma$-algebra that contains $\mathcal{F}$:
complements of elements of $\mathcal{F}$, countable unions of elements of
$\mathcal{F}$, countable unions of complements of elements of $\mathcal{F}$,
and so on --- iterating indefinitely. The generated $\sigma$-algebra is the
``closure'' of $\mathcal{F}$ under all $\sigma$-algebraic operations. In
practice, one specifies $\mathcal{F}$ as a convenient collection (e.g.,
intervals) and then works with the full $\sigma$-algebra by knowing which
operations are allowed.
\end{remark*}

\begin{definitionbox}{Definition (Borel $\sigma$-Algebra)}
\begin{definition}[Borel $\sigma$-Algebra]\label{def:borel}
Let $(X, \tau)$ be a topological space. The \emph{Borel $\sigma$-algebra}
on $X$, written $\mathcal{B}(X)$, is the $\sigma$-algebra generated by the
open sets:
\[
\mathcal{B}(X) \;:=\; \sigma(\tau).
\]
Members of $\mathcal{B}(X)$ are called \emph{Borel sets}.
\end{definition}
\end{definitionbox}

\begin{dependencies}
\begin{itemize}
  \item \hyperref[def:topology-defn]{Topology}
  \item \hyperref[def:generated-sigma]{Generated $\sigma$-Algebra}
\end{itemize}
\end{dependencies}

\begin{example*}[$F_\sigma$ and $G_\delta$ sets]
\label{ex:Fsigma-Gdelta}
In a topological space $X$:
\begin{itemize}
\item A \emph{$G_\delta$ set} is a countable intersection of open sets:
  $A = \bigcap_{n=1}^\infty U_n$ with each $U_n$ open.
\item An \emph{$F_\sigma$ set} is a countable union of closed sets:
  $A = \bigcup_{n=1}^\infty F_n$ with each $F_n$ closed.
\end{itemize}
These names come from German: $G$ for \emph{Gebiet} (open set), $\delta$ for
\emph{Durchschnitt} (intersection); $F$ for \emph{Ferme} (closed), $\sigma$
for \emph{Summe} (union). Both $G_\delta$ and $F_\sigma$ sets are Borel sets.
They are dual: $A$ is $G_\delta$ if and only if $A^c$ is $F_\sigma$.

These classes appear in analysis and measure theory as the natural ``next
level'' above open and closed sets. In $\mathbb{R}$: the set of continuity
points of a monotone function is a $G_\delta$ set; the rational numbers
$\mathbb{Q}$ are an $F_\sigma$ set; the irrational numbers $\mathbb{R}
\setminus \mathbb{Q}$ are a $G_\delta$ set.
\end{example*}

\begin{remark*}[Transition]
The four structures --- semiring, ring, algebra, $\sigma$-algebra --- and the
topology form a hierarchy of increasing closure strength. Measure theory
builds from the bottom (semiring, ring) upward; topology is its own branch
with a different closure pattern. The Borel $\sigma$-algebra bridges the two:
it is generated by the topology, so it captures all sets describable by
topological operations, and it is the standard ``default'' $\sigma$-algebra
in any analytic context. Whenever a later chapter says ``let $f$ be a
measurable function'' without specifying the $\sigma$-algebra, it almost
certainly means measurability with respect to the Borel $\sigma$-algebra.
\end{remark*}
